\section{Conclusion}\label{sec:conclusion}

The minimum integer count for quadratic graph approximation is governed
by a joint algebraic invariant: half the noncommutative rank of the
Hessian space is its exact logarithmic precision coefficient, for
arbitrary convex integer lifts and for binary linear lifts alike, although
their finite-accuracy counts can differ. At finite accuracy, contact
covariance gives a determinant benchmark and constructive bounds with
explicit dimension losses. Positive structure, effective input
coordinates and the actual error body can reduce those losses, while
correlated domains can prevent naive transfers.

For scalar convex graphs, chord geometry makes binary linear lifts
competitive with arbitrary convex integer lifts within two variables.
Certified integration, rational endpoint algorithms and a common index
compiler give an eleven-bit polynomial rational comparison for dense
convex polynomials, with separate constructions for the specified sparse
polynomial and power representations. Coupled outputs require
simultaneous control, which original-output or positive-polar curvature
bases provide; a shared basis across separable inputs preserves the
relevant packing lower bound. These results depend on convexity and the
stated error geometry. The exact product counts, positive-power
obstruction, nonconvex polynomial family and tilted-body example show
different ways in which those assumptions matter.

Several questions remain open. The finite covariance benchmark has an
order-$n\log n$ gap on an explicit family; a different efficiently
computable invariant might give smaller losses, subject to the proved
count-approximation hardness. Local and global smooth Hessian ranks can
differ, and the constant-rank theorem does not characterize rank-changing
scalar maps. For one input and arbitrarily many convex outputs with box
error, it is open whether a universal constant additive gap exists
(\cref{subsec:one-input-open}). Extending the separable vector construction to
mixed-coordinate nonlinear terms requires a new compact encoding
argument.

The encoding separation marks a further boundary: four binary variables
can require rational linear data exponentially long in the bit length of
the exponent, whereas second-order cone constraints achieve the same
graph accuracy with short data. Integer count, continuous size,
coefficient length, construction time and optimization time should
therefore be kept separate when these formulations are used or compared.
